\exercisegroup

\begin{enumerate}
\def\labelenumi{\arabic{enumi}.}
\tightlist
\item
  设 \(A, B, C\) 为逻辑理论 \(\mathfrak{G}\) 中的关系。证明以下关系是 \(\mathfrak{G}\) 中的定理：
\end{enumerate}

\[
\begin{array}{c} A \Longrightarrow (B \Longrightarrow A), \\ (A \Longrightarrow B) \Longrightarrow ((B \Longrightarrow C) \Longrightarrow (A \Longrightarrow C)), \\ A \Longrightarrow ((\text {非} A) \Longrightarrow B), \\ (A \text {或} B) \Longleftrightarrow ((A \Longrightarrow B) \Longrightarrow B), \\ (A \Longleftrightarrow B) \Longleftrightarrow ((A \text {且} B) \text {或} ((\text {非} A) \text {且} (\text {非} B))), \\ (A \Longleftrightarrow B) \Longleftrightarrow \text {非} ((\text {非} A) \Longleftrightarrow B), \\ (A \Longrightarrow (B \text {或} (\text {非} C))) \Longleftrightarrow ((C \text {且} A) \Longrightarrow B), \\ (A \Longrightarrow (B \text {或} C)) \Longleftrightarrow (B \text {或} (A \Longrightarrow C)), \\ (A \Longrightarrow B) \Longrightarrow ((A \Longrightarrow C) \Longrightarrow (A \Longrightarrow (B \text {且} C))), \\ (A \Longrightarrow C) \Longrightarrow ((B \Longrightarrow C) \Longrightarrow ((A \text {或} B) \Longrightarrow C)), \\ (A \Longrightarrow B) \Longrightarrow ((A \text {且} C) \Longrightarrow (B \text {且} C)), \\ (A \Longrightarrow B) \Longrightarrow ((A \text {或} C) \Longrightarrow (B \text {或} C)). \end{array}
\]

\begin{enumerate}
\def\labelenumi{\arabic{enumi}.}
\setcounter{enumi}{1}
\item
  设 \(A\) 为逻辑理论 \(\mathfrak{C}\) 中的一个关系。如果 \(A \Longleftrightarrow (\text{非 } A)\) 是 \(\mathfrak{C}\) 中的一个定理，那么 \(\mathfrak{C}\) 是矛盾的。
\item
  设 \(A_{1}, A_{2}, \ldots, A_{n}\) 为逻辑理论T中的关系。
\end{enumerate}

\begin{enumerate}
\def\labelenumi{(\alph{enumi})}
\item
  为了证明关系“ \(A_{1}\) 要么 \(A_{2}\) 要么 \ldots{} 要么 \(A_{n}\) ”在 \(\mathfrak{C}\) 中成立，只需证明 \(A_{n}\) 在通过将 \(\mathfrak{C}\) 公理不 \(A_{1}\) ，不 \(A_{2}\) , \ldots，不 \(A_{n-1}\) .
\item
  如果“ \(A_{1}\) 要么 \(A_{2}\) 要么 \ldots{} 要么 \(A_{n}\) ”是 \(\mathfrak{C}\) ，那么要证明一个定理 \(A\) 在 \(\mathfrak{C}\) 只需证明这些定理
\end{enumerate}

\[
A _ {1} \Longrightarrow A, \quad A _ {2} \Longrightarrow A, \quad \dots , \quad A _ {n} \Longrightarrow A.
\]

\begin{enumerate}
\def\labelenumi{\arabic{enumi}.}
\setcounter{enumi}{3}
\tightlist
\item
  设 A 和 B 是逻辑理论 G 中的关系。设 \(A|B\) 表示关系“(非 A) 或 (非 B)”。在 G 中证明以下定理：
\end{enumerate}

\[
\begin{array}{c} (\text {非} \boldsymbol {A}) \Longleftrightarrow (\boldsymbol {A} | \boldsymbol {A}), \\ (\boldsymbol {A} \text {或} \boldsymbol {B}) \Longleftrightarrow ((\boldsymbol {A} | \boldsymbol {A}) | (\boldsymbol {B} | \boldsymbol {B})), \\ (\boldsymbol {A} \text {且} \boldsymbol {B}) \Longleftrightarrow ((\boldsymbol {A} | \boldsymbol {B}) | (\boldsymbol {A} | \boldsymbol {B})), \\ (\boldsymbol {A} \Longrightarrow \boldsymbol {B}) \Longleftrightarrow (\boldsymbol {A} | (\boldsymbol {B} | \boldsymbol {B})). \end{array}
\]

\begin{enumerate}
\def\labelenumi{\arabic{enumi}.}
\setcounter{enumi}{4}
\tightlist
\item
  设 \(\mathfrak{G}\) 是具有显式公理 \(A_{1}, A_{2}, \ldots, A_{n}\) 的逻辑理论。则 \(A_{n}\) 独立于 \(\mathfrak{G}\) 的其他公理（ \(\S 2\) ，练习）当且仅当如下理论不是矛盾的：其符号和方案与 \(\mathfrak{G}\) 的相同，且其显式公理为 \(A_{1}, A_{2}, \ldots, A_{n-1}\) 与``不 \(A_{n}\) ''。
\end{enumerate}

